% Minimal two-column physics paper template.
% Compile with: pdflatex main.tex && bibtex main && pdflatex main.tex && pdflatex main.tex

\documentclass[twocolumn,aps,prd,superscriptaddress]{revtex4-2}

\usepackage{amsmath}
\usepackage{amssymb}
\usepackage{graphicx}
\usepackage{hyperref}
\usepackage[T1]{fontenc}

\begin{document}

\title{Paper Title}

\author{Author Name}
\email{author@example.com}
\affiliation{Department, University, City, Country}

\begin{abstract}
% 4--5 sentences following this structure:
% 1. Setting + motivation (1 sentence): physical system and why it matters.
% 2. What we do (1--2 sentences): method and scope. Begin with "We ...".
% 3. Main quantitative result (1 sentence): a specific number or comparison.
% 4. Broader implication (1 sentence): how this complements existing work.
\end{abstract}

\maketitle

% ---------------------------------------------------------------
\section{Introduction}
% ---------------------------------------------------------------

% Paragraph 1 (3--5 sentences): broad physical setting.

% Paragraph 2 (3--5 sentences): specific opportunity and prior work.

% Paragraph 3 (2--4 sentences): what this paper does.
% Begin with "In this work, we ..." or "In this paper, we show ...".

% Paragraph 4 (optional): paper structure.
% "In Section~\ref{sec:methods}, we ... We conclude in Section~\ref{sec:conclusion}."

% ---------------------------------------------------------------
\section{Methods}
\label{sec:methods}
% ---------------------------------------------------------------

% ---------------------------------------------------------------
\section{Results}
\label{sec:results}
% ---------------------------------------------------------------

% ---------------------------------------------------------------
\section{Conclusion}
\label{sec:conclusion}
% ---------------------------------------------------------------

% Paragraph 1: restate the main result in one sentence.

% Paragraph 2: broader implications.

% Paragraph 3: outlook -- two or three natural follow-up directions.

\bibliography{references}

\end{document}
